\documentclass[12pt]{article}

%% ── Packages ────────────────────────────────────────────────────────────────
\usepackage{amsmath,amssymb,amsthm}
\usepackage{algorithm}
\usepackage{algpseudocode}
\usepackage{booktabs}
\usepackage{graphicx}
\usepackage{hyperref}
\usepackage{xcolor}
\usepackage{microtype}
\usepackage{geometry}
\usepackage{natbib}
\usepackage{setspace}
\usepackage{caption}
\usepackage{subcaption}
\usepackage{multirow}

\geometry{margin=1in}
\onehalfspacing

%% ── Theorem environments ────────────────────────────────────────────────────
\newtheorem{definition}{Definition}
\newtheorem{proposition}{Proposition}
\newtheorem{remark}{Remark}

%% ── Symbol shorthands ───────────────────────────────────────────────────────
\newcommand{\cS}{\mathcal{S}}       % raw state space
\newcommand{\cA}{\mathcal{A}}       % action space
\newcommand{\cC}{\mathcal{C}}       % symbolic condition space
\newcommand{\cR}{\mathcal{R}}       % rule bank
\newcommand{\cG}{\mathcal{G}}       % grammar / rule set
% \Phi is a standard LaTeX math symbol; no redefinition needed
\newcommand{\piN}{\pi_{\theta}}     % neural policy
\newcommand{\Jac}{\mathrm{Jac}}     % Jaccard
\newcommand{\PA}{\mathrm{PA}}       % provenance audit
\newcommand{\deltaDelta}{\Delta}    % paired return difference

%% ────────────────────────────────────────────────────────────────────────────
\begin{document}

\title{\textbf{Provenance-Preserving Rule-Level Fusion of\\
Reinforcement-Learning Policies}\\[4pt]
\large Explicit Conflicts, Auditable Traces, and the Gap Between\\
Symbolic Overlap and Behavioral Consensus}

\author{Anonymous Submission}
\date{}
\maketitle

%% ── Abstract ────────────────────────────────────────────────────────────────
\begin{abstract}
Reinforcement-learning (RL) systems are difficult to audit because
trajectory replay, decision explanation, and multi-agent knowledge composition
are three distinct properties that are often conflated under the single label
\emph{auditability}.
We present a framework that separates these properties and evaluates each
empirically on two standard benchmark environments.

Our central mechanism is a \emph{provenance-preserving rule merger}:
given independently trained agents whose state observations are discretised
through a shared frozen symboliser, the merger records, for every decision,
the candidate rules from each source, the arbitration outcome, the fallback
agent when no rule applies, and a chained hash linking the decision to the
original training record.

We conduct experiments with eight REINFORCE seeds on
\textsc{CartPole-v1} and \textsc{Acrobot-v1}.
Under a shared symboliser, eight-seed pairwise rule-set
Jaccard is $0.557$ (mean) for CartPole and $0.096$ for Acrobot,
revealing that symbolic overlap is highly task-dependent.
A direct held-out action-agreement experiment (\emph{Experiment~B})
shows that the two CartPole seeds with the highest pairwise Jaccard ($0.583$)
agree on only $51.4\%$ of fresh raw observations
(95\% CI $[50.5\%, 52.3\%]$), indistinguishable from the binary chance baseline.
Symbolic overlap therefore cannot serve as a proxy for behavioral consensus.

Grammar-based fusion achieves a statistically significant return improvement
over the best single agent on both tasks
(CartPole: $+42.3$ points, 95\% CI $[24.0, 61.9]$;
Acrobot: $+56.6$, 95\% CI $[39.3, 75.2]$),
yet confidence-ranked arbitration performs substantially worse than
random rule selection on Acrobot ($+237$--$249$ units for random,
CI non-overlapping with zero),
indicating that arbitration mechanism design is the primary performance driver,
not the size of the shared rule core.
All fused traces pass exact environment replay and source-rule
provenance verification.
These results support bounded claims: provenance-linked symbolic
rule composition makes selected decisions auditable under controlled
conditions; they do not establish general RL training transparency
or a performance advantage over neural ensemble baselines.
\end{abstract}

\tableofcontents
\newpage

%% ════════════════════════════════════════════════════════════════════════════
\section{Introduction}
\label{sec:intro}
%% ════════════════════════════════════════════════════════════════════════════

%% ── 1.1 Scenario ────────────────────────────────────────────────────────────
Consider an operator who wishes to deploy a policy composed from several
independently trained RL agents and must later explain why a particular action
was taken at a particular state.
Existing approaches offer partial answers.
Neural policy ensembles average logits from multiple networks but produce
no record of which network drove which decision.
Post-hoc saliency methods generate explanations after the fact with no
guarantee that the explanation corresponds to the computation performed.
Programmatic RL~\citep{verma2018pirl} encodes policies as human-readable
programs, but requires re-expressing the policy in a pre-designed
domain-specific language.

%% ── 1.2 Limitations of existing approaches ──────────────────────────────────
Three limitations motivate the present work.
\emph{First}, trajectory replay and decision explanation are not the same:
a system can log every action deterministically and still provide no
interpretable rationale for each choice.
\emph{Second}, symbolic overlap between rule sets from different agents
does not imply behavioral consensus on the same raw states, as we demonstrate
empirically.
\emph{Third}, multi-agent policy composition methods such as Generalized
Policy Improvement~(GPI)~\citep{barreto2017successor} require value-function
estimation whose quality degrades under standard function approximation, and
provide no mechanism for recording the \emph{source} of each merged decision.

%% ── 1.3 Problem and objective ───────────────────────────────────────────────
We address the following question: \emph{can a rule-level merger make the
provenance of every fused decision explicitly verifiable, expose agent
disagreements and uncovered states, and do so without requiring a
pre-designed symbolic language?}
The objective is audit visibility, not necessarily higher return.

%% ── 1.4 Challenges ──────────────────────────────────────────────────────────
The challenge is that learned symbolic rules carry three sources of opacity:
(i)~the discretisation of the continuous state space introduces vocabulary
mismatch when different agents use independently fitted quantisers;
(ii)~two agents may assign the same discrete symbol to different raw-state
regions, making their rules syntactically identical but semantically
divergent;
(iii)~when rules conflict on a state, the choice of arbitration order
has empirically large effects on performance that cannot be predicted from
rule statistics alone.

%% ── 1.5 Method components ───────────────────────────────────────────────────
We address these challenges through three components.
A \emph{shared frozen symboliser} eliminates vocabulary mismatch by fitting
a single discretisation from pooled random-policy rollouts before any agent
trains.
A \emph{provenance-aware merger} records, at every decision step, all
candidate rules, the arbitration path taken, and the source agent used for
fallback, then binds the decision record into a hash chain that allows
post-hoc tampering detection.
A \emph{replay-and-provenance auditor} independently re-executes the
environment and resolves every rule reference back to the original training
record, confirming that the logged and actual sequences match and that each
cited rule existed with its stated statistics at the time of the decision.

%% ── 1.6 Contributions ───────────────────────────────────────────────────────
Our contributions are:
\begin{enumerate}
\item A provenance-preserving rule merger for independently trained RL agents
  that makes conflicts, blind spots, arbitration paths, and fallback sources
  explicit and verifiable at decision time
  (Sections~\ref{sec:method:merger}--\ref{sec:method:auditor}).
\item An empirical decomposition of symbolic overlap into a syntactic
  component (Jaccard on rule sets) and a behavioral component (held-out
  action agreement), showing that the former substantially overstates the
  latter in the tested settings
  (Section~\ref{sec:results:exptB}).
\item A multi-seed, two-environment evaluation under a shared symboliser
  revealing that rule-set overlap is task-dependent and that arbitration
  mechanism design dominates composition performance
  (Section~\ref{sec:results:composition}).
\item A causal-symbol intervention result in a separate designed communication
  game, supporting the claim that learned discrete symbols can carry
  task-relevant information, with explicit scope boundaries
  (Section~\ref{sec:results:aim}).
\end{enumerate}

%% ════════════════════════════════════════════════════════════════════════════
\section{Background and Related Work}
\label{sec:related}
%% ════════════════════════════════════════════════════════════════════════════

\paragraph{Interpretable and programmatic RL.}
Programmatically Interpretable Reinforcement Learning
(PIRL,~\citet{verma2018pirl}) synthesises policies in a domain-specific
language, supporting symbolic inspection and, in some formulations, formal
verification.
The present work asks a different question: rather than re-expressing a
trained policy in a pre-designed language, can we induce a symbolic rule
bank from the agent's own trajectory and use it to compose multiple agents
with traceable provenance?
The approaches are complementary: PIRL trades neural flexibility for
full policy transparency, whereas our merger preserves the neural policy
as a fallback while exposing the symbolic layer's coverage and conflicts.

\paragraph{Policy composition via successor features.}
\citet{barreto2017successor} introduce successor features (SFs) and show
that Generalised Policy Improvement (GPI)~\citep{barreto2017gpi} can
compose policies by taking the argmax of per-source value estimates.
\citet{thakoor2022geometric} extend this to geometric composition.
These methods compose knowledge at the level of value functions and require
a critic of sufficient quality; in our experiments, a fitted
Q-function evaluator failed pre-specified action-quality checks on both
environments and was excluded from primary comparisons
(Section~\ref{sec:results:gpi}).
More fundamentally, value-function composition does not record \emph{which}
source agent contributed each individual decision.

\paragraph{Emergent communication and symbol causality.}
Work on emergent communication demonstrates that agents can develop
referential protocols~\citep{lazaridou2018emergent}, but cautions that
message--action correlation alone does not establish causal influence;
intervention is required~\citep{lowe2019pitfalls}.
We apply a matched-support intervention in a designed private-target
game and report per-seed causal effect estimates
(Section~\ref{sec:results:aim}).

\paragraph{Reproducibility in deep RL.}
Exact trajectory replay depends on seed control and
software/environment state~\citep{nagarajan2018deterministic}.
We design the audit to distinguish \emph{exact replay} of logged actions
from \emph{fresh generation} by a trained policy,
and limit replay claims to the tested hardware and software configuration.

\paragraph{Sequence grammar induction.}
SEQUITUR~\citep{nevill1997sequitur} induces a hierarchical context-free
grammar from a symbol sequence by replacing repeated bigrams with
non-terminal symbols.
We use a two-part minimum-description-length (MDL) criterion,
$L(\cG) + L(\tau \mid \cG)$, to evaluate whether the induced grammar
compresses a trajectory.
The grammar archive alone is not a valid compression benchmark against
raw gzip; only the combined two-part length is informative.

%% ════════════════════════════════════════════════════════════════════════════
\section{Problem Formulation}
\label{sec:problem}
%% ════════════════════════════════════════════════════════════════════════════

Let $(\cS, \cA, P, r, \gamma)$ be a Markov decision process with
state space $\cS \subseteq \mathbb{R}^d$, finite action space
$\cA = \{0,\ldots,|\cA|-1\}$, transition kernel $P$, reward $r$, and
discount $\gamma$.
Let $\piN^{(k)} : \cS \to \Delta(\cA)$ denote the neural policy of the
$k$-th independently trained agent, $k = 1,\ldots,K$.

\begin{definition}[Shared symboliser]
\label{def:symboliser}
A shared symboliser $\Phi : \cS \to \cC$ maps each raw observation to a
discrete condition symbol.
$\Phi$ is fitted from pooled random-policy rollouts before any agent trains
and is frozen thereafter, so that every agent operates in the same symbolic
vocabulary $\cC$.
In our experiments, $\Phi$ partitions each state dimension into $B = 4$ bins
using empirical quantile boundaries; $|\cC| = B^d$.
\end{definition}

\begin{definition}[State-action rule bank]
\label{def:rulebank}
Given agent $k$ and symboliser $\Phi$, the rule bank
$\cR^{(k)} \subseteq \cC \times \cA$ is the set of pairs $(c, a)$ such that
the observed frequency of action $a$ given condition $\Phi(s) = c$ meets
a minimum support threshold $n_{\min}$ and confidence threshold
$\kappa_{\min}$:
\begin{equation}
\label{eq:rulebank}
\cR^{(k)} = \bigl\{(c, a) \mid
  \mathrm{count}(c, a) \ge n_{\min},\;
  \tfrac{\mathrm{count}(c,a)}{\mathrm{count}(c)} \ge \kappa_{\min}\bigr\}.
\end{equation}
The rule bank is distinct from the neural policy: it is a sparse lookup
table, not a function approximator.
\end{definition}

\begin{definition}[Jaccard rule overlap]
\label{def:jaccard}
For two rule banks $\cR^{(i)}$ and $\cR^{(j)}$, the Jaccard coefficient is
\begin{equation}
\label{eq:jaccard}
\Jac\bigl(\cR^{(i)}, \cR^{(j)}\bigr)
  = \frac{|\cR^{(i)} \cap \cR^{(j)}|}{|\cR^{(i)} \cup \cR^{(j)}|}.
\end{equation}
Note that $\Jac$ measures syntactic identity of $(c,a)$ pairs under the
shared symboliser; it does not measure agreement on raw states outside
the training distribution.
\end{definition}

\begin{definition}[Behavioral agreement]
\label{def:bahavior_agree}
The held-out behavioral agreement between agents $i$ and $j$ is
\begin{equation}
\label{eq:behavioral_agree}
\mathrm{BA}(i, j)
  = \mathbb{E}_{s \sim \mu_{\mathrm{held-out}}}
    \bigl[\mathbf{1}\bigl\{\arg\max_a \piN^{(i)}(a \mid s)
          = \arg\max_a \piN^{(j)}(a \mid s)\bigr\}\bigr],
\end{equation}
where $\mu_{\mathrm{held-out}}$ is a fresh state distribution obtained by
rolling out each agent's deterministic policy on independent reset seeds
not used during training or primary evaluation.
\end{definition}

Definitions~\ref{def:jaccard} and~\ref{def:bahavior_agree} measure
different quantities; the relationship between them is an empirical question,
not a logical implication.

%% ════════════════════════════════════════════════════════════════════════════
\section{Method}
\label{sec:method}
%% ════════════════════════════════════════════════════════════════════════════

\subsection{Shared Symboliser and Rule Induction}
\label{sec:method:symboliser}

For each environment, we pool $N_{\mathrm{cal}}$ episodes of a
random policy across all calibration seeds and fit per-dimension empirical
quantile boundaries with $B$ bins.
The resulting boundaries are frozen into a canonical artefact whose
SHA-256 hash is recorded.
Every subsequent training arm, audit arm, and evaluation uses the identical
symboliser, ensuring that the symbolic vocabulary is common across agents
and that condition tokens from different seeds refer to the same raw-state
regions.

Rule banks are induced from the trajectory of the audit-only training arm,
which observes state--action pairs without modifying the neural policy update.
We verify that the baseline and audit-only arms produce identical neural
weights and transition records for every seed and environment.

\subsection{Provenance-Aware Rule Merger}
\label{sec:method:merger}

Algorithm~\ref{alg:merger} describes the per-step decision procedure of the
merged policy $\pi_{\mathrm{merge}}$.

\begin{algorithm}[t]
\caption{Provenance-preserving rule merger}
\label{alg:merger}
\begin{algorithmic}[1]
\Require Raw state $s \in \cS$;
         rule banks $\{\cR^{(k)}\}_{k=1}^{K}$;
         symboliser $\Phi$;
         fallback neural policies $\{\piN^{(k)}\}$;
         fallback selector $k^* = \arg\max_k \bar{G}^{(k)}_{\mathrm{train}}$
\Ensure  Action $a \in \cA$; provenance record $\rho$
\State $c \leftarrow \Phi(s)$
\State $\mathcal{M} \leftarrow \{(k, a_k) \mid (c, a_k) \in \cR^{(k)}\}$
  \Comment{Collect matching rules from all sources}
\If{$|\mathcal{M}| = 0$}
  \State $a \leftarrow \arg\max_a \piN^{(k^*)}(a \mid s)$
  \Comment{Blind-spot fallback; no rule applies}
  \State $\rho.\mathrm{type} \leftarrow \textsc{blind-spot}$
\ElsIf{all $a_k$ in $\mathcal{M}$ are equal}
  \State $a \leftarrow a_k$ for any $k \in \mathcal{M}$
  \State $\rho.\mathrm{type} \leftarrow \textsc{unanimous}$
\Else
  \State $\rho.\mathrm{type} \leftarrow \textsc{conflict}$
  \State Rank candidates by $(\kappa^{(k)}_c, n^{(k)}_c, \bar{r}^{(k)}_c)$
    descending \label{alg:line:rank}
  \State $a \leftarrow$ action from top-ranked candidate
\EndIf
\State $\rho \leftarrow \bigl(c,\, a,\, \mathcal{M},\,
  \mathrm{hash}(\text{prev. record})\bigr)$
\State \Return $a$, $\rho$
\end{algorithmic}
\end{algorithm}

In Algorithm~\ref{alg:merger}, line~\ref{alg:line:rank}, candidates are
ranked lexicographically by per-condition confidence $\kappa^{(k)}_c$,
support $n^{(k)}_c$, and mean observed reward $\bar{r}^{(k)}_c$, in that
order.
Exact ties use the fallback neural policy.
The fallback agent $k^*$ is selected by training-episode mean return only;
held-out evaluation returns are never consulted to avoid contamination.

Each provenance record $\rho$ stores:
(i)~the symbolic condition $c$,
(ii)~all candidate rule identifiers and their source-step hash digests,
(iii)~the selected source(s) or the fallback identity,
(iv)~a chained record hash $h_t = H(\rho_t \| h_{t-1})$ where $H$ is SHA-256.

\subsection{Replay and Provenance Auditor}
\label{sec:method:auditor}

The auditor performs two independent verification passes.

\paragraph{Environment replay.}
For each logged episode, the auditor resets the environment with the recorded
seed and replays the logged action sequence, comparing the observed next
state, reward, termination flag, and truncation flag to the logged values
at every step.
A mismatch at any field constitutes a replay failure.

\paragraph{Provenance resolution.}
For each constrained decision, the auditor resolves the cited rule identifier
to the source rule bank, verifies that the rule's support and confidence meet
the declared thresholds, checks that the cited source-step hash digests
appear in the training ledger, and confirms that the emitted action matches
the rule's action.
The hash chain is re-computed from the first record and compared to the
stored chain; any discrepancy indicates post-hoc modification.

\begin{definition}[Replay-and-provenance audit]
\label{def:audit}
A fused trace $\tau = (\rho_1, \ldots, \rho_T)$ passes the
replay-and-provenance audit $\PA(\tau) = \mathrm{true}$ if and only if:
\begin{align}
\label{eq:audit}
&\forall t:\; s'_{t,\mathrm{env}} = s'_{t,\mathrm{log}},\quad
  r_{t,\mathrm{env}} = r_{t,\mathrm{log}},\quad
  d_{t,\mathrm{env}} = d_{t,\mathrm{log}},\\
&\forall t \text{ with rule decision}:\;
  (c_t, a_t) \in \cR^{(k_t)},\quad
  n^{(k_t)}_{c_t} \ge n_{\min},\quad
  \kappa^{(k_t)}_{c_t} \ge \kappa_{\min},\nonumber\\
&\forall t:\; h_t = H(\rho_t \| h_{t-1}),\nonumber
\end{align}
where subscripts $\mathrm{env}$ and $\mathrm{log}$ denote the
independently replayed and originally logged values, respectively.
\end{definition}

\subsection{MDL Trajectory Coding}
\label{sec:method:mdl}

To characterise the descriptive complexity of a logged trajectory, we
adopt the two-part minimum-description-length (MDL) criterion.
Let $\cG_\tau$ be a context-free grammar induced from trajectory $\tau$
by the SEQUITUR algorithm~\citep{nevill1997sequitur}.
Define the two-part description length
\begin{equation}
\label{eq:mdl}
L_2(\tau) = L(\cG_\tau) + L(\tau \mid \cG_\tau),
\end{equation}
where $L(\cG_\tau)$ is the encoded grammar size and
$L(\tau \mid \cG_\tau)$ is the length of the trace encoded under the grammar.
The compression ratio $\rho_{\mathrm{MDL}} = L_2(\tau) / L(\tau)$
compares the two-part code to the raw trace length under a reference code.
Note that $\cG_\tau$ is a \emph{lossless trace grammar} that reconstructs
the observed sequence exactly; it is not the state-action rule bank
$\cR^{(k)}$ and does not claim to explain the neural policy.

\subsection{Experimental Protocol}
\label{sec:method:protocol}

\paragraph{Agents and environments.}
We train eight REINFORCE agents (seeds 11, 29, 43, 71, 101, 149, 211, 307)
on \textsc{CartPole-v1} ($d=4$, $|\cA|=2$) and \textsc{Acrobot-v1}
($d=6$, $|\cA|=3$) for 300 episodes each.

\paragraph{Shared symboliser.}
For each environment, we pool $N_{\mathrm{cal}} = 24$ random-policy episodes
per calibration seed, fit $B = 4$ empirical quantile bins per state
dimension, and freeze the symboliser.
Rule admission requires minimum support $n_{\min} = 8$ and confidence
$\kappa_{\min} = 0.70$.

\paragraph{Arms.}
Each seed trains three parallel arms: (i)~a \emph{baseline} arm with no
observer, (ii)~an \emph{audit-only} arm whose observer records symbolic
transitions without modifying gradient updates, and (iii)~a
\emph{grammar-constrained} arm that applies the current rule bank as a
soft mask during training.
Baseline and audit-only arms are verified to produce identical network
weights and transition records for every seed.

\paragraph{Composition evaluation.}
Eight source actors and the grammar-fusion policy are evaluated on $N_{\mathrm{eval}} = 100$
common reset seeds
(formula: $2 \times 10^7 + \text{episode index}$, disjoint from training
and Experiment B resets).
The fallback agent $k^*$ is the seed with the highest training mean return,
selected without reference to evaluation performance.

\paragraph{Estimands and predictions.}
Before observing evaluation results, we register four predictions:
\begin{enumerate}
\item[P1] Grammar fusion achieves a non-negative return difference versus
  the best single agent on CartPole (positive because the merger can always
  fall back to the best single agent).
\item[P2] Grammar fusion achieves a non-negative return difference versus
  the best single agent on Acrobot.
\item[P3] Held-out behavioral agreement (Def.~\ref{def:bahavior_agree})
  for the two seeds with highest pairwise Jaccard exceeds $85\%$.
\item[P4] The two-part MDL ratio $\rho_{\mathrm{MDL}}$ is less than 1.0
  for all trajectory prefix lengths tested (i.e., the grammar represents
  a genuine compression of the trace).
\end{enumerate}
All evaluation uses matched-seed bootstrap resampling
($10^4$ resamples) for 95\% confidence intervals.

%% ════════════════════════════════════════════════════════════════════════════
\section{Results}
\label{sec:results}
%% ════════════════════════════════════════════════════════════════════════════

\subsection{Audit Integrity}
\label{sec:results:audit}

\paragraph{Replay.}
The auditor verified exact environment replay for
1{,}122{,}403 CartPole transitions and 1{,}742{,}570 Acrobot transitions
across all 48 training arms (8 seeds $\times$ 2 environments $\times$ 3 arms).
No replay mismatch was observed.

\paragraph{Hash-chain and provenance.}
All 48 hash chains passed re-computation.
All baseline--audit-only paired arms were confirmed to be exactly equivalent
(identical network weights and transition records) for every seed on both
environments.
Provenance resolution succeeded for all 290{,}052 constrained decisions in
CartPole and all 697{,}817 in Acrobot:
every cited rule reference resolved to a source rule bank meeting the
declared thresholds, and every source-step hash appeared in the corresponding
training ledger.
A tamper probe that modified the reward in the first record of a trace copy
produced immediate hash-chain failure at step 1.

\paragraph{Verdict.}
$\PA(\tau) = \mathrm{true}$ for all evaluated traces.
The audit confirms logger consistency and post-hoc tamper detectability
on the tested hardware and software configuration.
It does not establish that the logger was honest during training, that
the record generalises to other hardware, or that the symbolic rules
causally explain the neural policy's computation.

\subsection{Rule-Set Overlap Under the Shared Symboliser}
\label{sec:results:overlap}

\paragraph{CartPole.}
Under the shared symboliser, the eight-seed mean pairwise Jaccard is
$0.557$ (median $0.565$).
The all-seed unanimous-action intersection contains 38 condition symbols;
no condition in this intersection exhibits a cross-seed action conflict.
Per-seed rule counts range from 100 to 145 (mean 115).
The seed-11 vs.\ seed-29 pairwise Jaccard is $0.583$, with 77 shared
conditions and zero action conflicts.

\paragraph{Acrobot.}
The eight-seed mean pairwise Jaccard is $0.096$ (median $0.059$).
The all-seed intersection contains only 2 condition symbols, both of
which exhibit action conflicts across seeds.
The all-seed unanimous-action core contains zero rules.
The seed-11 vs.\ seed-29 pairwise Jaccard is $0.050$, with 21 shared
conditions of which 77\% involve action conflicts.

\paragraph{Task dependence.}
CartPole's low-dimensional, nearly deterministic balance task produces a
large symbolic core under the shared symboliser; Acrobot's higher-dimensional
swing-up task produces almost no stable core across seeds.
This contrast indicates that symbolic rule overlap is governed by task
structure, not by the framework itself.

Table~\ref{tab:overlap} summarises the overlap statistics.

\begin{table}[t]
\centering
\caption{Rule-set overlap statistics under the shared frozen symboliser.
  All Jaccard values use $\kappa_{\min}=0.70$, $n_{\min}=8$.}
\label{tab:overlap}
\begin{tabular}{lcc}
\toprule
Statistic & CartPole-v1 & Acrobot-v1 \\
\midrule
8-seed mean pairwise Jaccard          & 0.557 & 0.096 \\
8-seed median pairwise Jaccard        & 0.565 & 0.059 \\
Seed-11 vs.\ seed-29 Jaccard          & 0.583 & 0.050 \\
All-seed unanimous-action core (rules)& 38    & 0     \\
Conditions with cross-seed conflicts  & 0     & 77/98 \\
Per-seed rule count (mean)            & 115   & 193   \\
\bottomrule
\end{tabular}
\end{table}

\subsection{Held-Out Behavioral Agreement (Experiment~B)}
\label{sec:results:exptB}

\paragraph{Design.}
\emph{Prediction P3} stated that held-out behavioral agreement between
the highest-Jaccard seed pair would exceed $85\%$.

We evaluated seeds 11 and 29 on CartPole-v1, the pair with the highest
pairwise Jaccard ($0.583$).
Each actor's deterministic policy was rolled out for 100 fresh episodes
on reset seeds in the range $[4\times10^7, 4\times10^7+99]$ (seed 11)
and $[4.1\times10^7, 4.1\times10^7+99]$ (seed 29), entirely disjoint from
training and primary evaluation resets.
We sampled 5{,}000 raw states without replacement from each actor's
state occupancy, giving $n = 10{,}000$ total.
Both actors were applied to every raw state; the episode-cluster bootstrap
corrects for intra-episode correlation.

\paragraph{Result.}
The pooled agreement rate was $51.4\%$
(episode-cluster balanced: $51.4\%$;
95\% bootstrap CI $[50.5\%, 52.3\%]$).
For the sub-sample drawn from seed-11 occupancy, agreement was $50.6\%$;
for the sub-sample from seed-29 occupancy, it was $52.2\%$.

The action histogram reveals the mechanism:
on seed-29 occupancy states, actor 29 splits actions roughly evenly
(0: $52\%$, 1: $48\%$), while actor 11 selects action 0 on $99.9\%$
of the same states.
This indicates that the two agents occupy largely non-overlapping regions
of the raw state space and are near-deterministic in each other's occupancy.
The polar disagreement at the occupancy boundary produces a pooled rate
close to the binary chance baseline of $50\%$.

\paragraph{Verdict.}
\emph{Prediction P3 is contradicted.}
Despite a Jaccard of $0.583$, the held-out behavioral agreement is
indistinguishable from chance.
Symbolic rule overlap, as measured by Jaccard under the shared symboliser,
is a product of overlapping training-state occupancies in the symbolic
vocabulary; it does not measure neural-policy consensus on arbitrary
held-out states.

\subsection{Composition Evaluation}
\label{sec:results:composition}

Table~\ref{tab:composition} reports mean episode return and paired
bootstrap differences versus the best single agent.

\begin{table}[t]
\centering
\caption{Composition evaluation on 100 common reset seeds.
  Paired delta and 95\% CI are against the best single agent.
  CartPole best-single seed: 43; Acrobot best-single seed: 211.}
\label{tab:composition}
\begin{tabular}{lcccc}
\toprule
Policy & \multicolumn{2}{c}{CartPole-v1} & \multicolumn{2}{c}{Acrobot-v1} \\
\cmidrule(lr){2-3}\cmidrule(lr){4-5}
       & Mean return & $\Delta$ (95\% CI) & Mean return & $\Delta$ (95\% CI) \\
\midrule
Best single actor        & 163.6  & ---            & $-500.0$ & --- \\
Actor-logit ensemble     & 500.0  & $+336.4$ $[326.1, 346.2]$ & $-500.0$ & $0.0$ \\
Grammar fusion           & 205.9  & $+42.3$ $[24.0, 61.9]$    & $-443.4$ & $+56.6$ $[39.3, 75.2]$ \\
MC-GPI (diagnostic only) & 9.4    & $-154.3$ $[-164.4,-144.6]$ & $-500.0$ & $0.0$ \\
\bottomrule
\end{tabular}
\end{table}

\paragraph{CartPole.}
\emph{Prediction P1 is supported.}
Grammar fusion yields a mean return of $205.9$ versus the best single
agent's $163.6$ ($\Delta = +42.3$, 95\% CI $[24.0, 61.9]$).
The actor-logit ensemble reaches the $500$-point ceiling, so CartPole
does not provide a meaningful performance comparison between neural
ensembles and grammar fusion.
The grammar-fusion improvement over the best single agent is
statistically significant, though the scale is saturated for the ensemble.

\paragraph{Acrobot.}
\emph{Prediction P2 is supported.}
Grammar fusion achieves mean return $-443.4$ versus $-500.0$ for both
the best single agent and the actor-logit ensemble
($\Delta = +56.6$, 95\% CI $[39.3, 75.2]$).
Acrobot's Acrobot's median return remains $-500$ (floor); the mean improvement
reflects 39 of 100 episodes with early goal termination (return $> -500$),
which the single-agent and ensemble policies could not achieve.

\paragraph{Coverage and conflicts on Acrobot.}
The fused policy was rule-covered on $95.7\%$ of steps and encountered
conflicts on $89.2\%$ of covered steps.
With $4.3\%$ blind spots, the composition performance is dominated by
conflict arbitration, not by uncovered-state fallback.

\paragraph{Arbitration ablations on Acrobot.}
Table~\ref{tab:arbitration} reports matched-seed ablations of the
arbitration mechanism.

\begin{table}[t]
\centering
\caption{Arbitration mechanism ablations on Acrobot-v1 (100 episodes,
  matched reset seeds). Random-rule arbitration uses 5 independent RNG
  streams; values are the range.
  Uniform-random-action control confirms the random-rule effect is not
  explained by unconditioned action noise.}
\label{tab:arbitration}
\begin{tabular}{lcc}
\toprule
Arbitration variant & Mean return & $\Delta$ vs.\ grammar fusion \\
\midrule
Grammar fusion (confidence/support/reward)  & $-443.4$ & --- \\
Confidence only                             & $-444.5$ & $-1.1$ \\
Support first                               & $-500.0$ & $-56.6$ \\
Mean-reward only                            & $-500.0$ & $-56.6$ \\
Defer conflicts to actor ensemble           & $-500.0$ & $-56.6$ \\
Random rule selection (5 streams)           & $-194.1$ to $-206.2$ & $+237$--$+249$ \\
\midrule
Uniform-random-action control (5 streams)   & $-496.4$ to $-499.9$ & $-53$ to $-57$ \\
\bottomrule
\end{tabular}
\end{table}

Random rule selection outperforms confidence-ranked arbitration by
$237$--$249$ return units (all five 95\% CIs exclude zero).
The uniform-random-action control averages $-498.7$, confirming that
the random-rule advantage is not explained by unconditioned action noise.
The mechanism underlying this reversal is not identified by the current
experiment; it requires follow-up analysis of state-occupancy distributions
and rule-source frequencies under each arbitration variant.

\paragraph{Verdict on arbitration.}
The confidence-ranked arbitration order is not an established improvement
over random rule selection on Acrobot.
This is a primary negative finding and affects the interpretation of
all composition results.

\subsection{MC-GPI and FQE Diagnostics}
\label{sec:results:gpi}

Fitted Q-evaluation (FQE) was used to fit per-source action-value
estimates for GPI action selection.
Pre-specified quality checks required that (i)~FQE actions agree with
the actor-logit ensemble on at least $50\%$ of held-out states, and
(ii)~value estimates are monotonically higher for held-out episodes
with higher observed return.

CartPole FQE-GPI agreed with the actor-logit ensemble on $35.3\%$ of
$31{,}563$ held-out states, below the pre-specified threshold.
Acrobot FQE-GPI selected action 0 on every state in a $50{,}000$-state
diagnostic set.
The Pearson correlation between source critics' Q-values on shared
held-out states was $0.65$, yet their greedy actions disagreed sharply.

Both checks failed.
FQE-GPI is reported in Table~\ref{tab:composition} as a diagnostic
reference only and is excluded from the primary composition comparison.

\subsection{AIM Symbol Causality}
\label{sec:results:aim}

As a separate experiment, we evaluate whether discrete learned symbols
in a private-target coordination game carry causal influence over
a receiver agent's action.
Six independent seeds of an AIM-based sender--receiver pair were trained
for $10^5$ episodes each.
A matched-support intervention replaced the information-bearing symbol
at the active token position with tokens from even and odd index sets
while holding the other token and receiver context constant.

\paragraph{Results.}
Across six seeds, sampled receiver accuracy was $0.999$ (mean);
shuffled-message accuracy was $0.503$ (mean);
no-message baseline was $0.500$ (mean).
The matched-token causal effect $\Delta P(B \mid C)$ ranged from
$0.997$ to $1.000$ across seeds (mean $0.998$).
Mutual information of the active token was $0.98$--$0.99$ bits across seeds.

\paragraph{Scope.}
The active token position varied across seeds (position~1 for four seeds,
position~0 for two), indicating that the network learns to use an available
channel without fixing its location.
These results support a causal role for learned discrete symbols in
this designed private-target game.
They do not establish symbol interpretability in arbitrary RL control
tasks, nor do they address the CartPole or Acrobot composition results.

\subsection{MDL Trajectory Coding}
\label{sec:results:mdl}

The two-part MDL ratio $\rho_{\mathrm{MDL}}$ (Equation~\ref{eq:mdl})
was computed for six successive 50-episode prefixes of a single
CartPole audit-only trajectory.
All six prefixes reconstructed exactly (lossless).
The ratio decreased from $0.513$ at prefix 1 to $0.575$ at prefixes 5--6,
where it stabilised.
\emph{Prediction P4 is supported:} $\rho_{\mathrm{MDL}} < 1.0$ at all
prefix lengths, confirming genuine two-part compression.

\begin{remark}
The grammar archive alone is larger than a raw gzip of the same trace
(grammar gzip: 13.4~MB; raw trace gzip: 12.3~MB).
This is expected: the lossless grammar is not optimised for archive
compression and must be evaluated under the declared two-part code.
The MDL ratio below 1.0 indicates that the grammar plus the
grammar-encoded trace together are shorter than the raw trace under
the same reference code.
\end{remark}

%% ════════════════════════════════════════════════════════════════════════════
\section{Discussion}
\label{sec:discussion}
%% ════════════════════════════════════════════════════════════════════════════

\paragraph{Why does symbolic overlap not imply behavioral consensus?}
The Jaccard coefficient under the shared symboliser measures whether
two agents assigned the same action to the same coarse symbolic cell
during their respective training trajectories.
When two agents' training-state occupancies overlap in the symbolic vocabulary
but differ substantially in the raw continuous space, they may map
to the same cell $c$ from different raw regions and yet select
opposite actions on held-out raw states from each other's occupancy.
Experiment~B reveals this mechanism: seed~29's actor selects action~0
on $99.9\%$ of seed-11 occupancy states, whereas seed~11 splits evenly
on those states.
The symbolic agreement on training conditions coexists with near-maximal
disagreement on held-out raw states from the partner's occupancy.
This is not a failure of the framework; it is a fundamental property of
coarse discretisation that the framework now makes measurable.

\paragraph{Why does random arbitration outperform confidence-ranked
arbitration on Acrobot?}
This is the most unexpected finding and currently lacks a mechanistic
explanation.
Three candidate explanations merit future investigation.
\emph{State-occupancy mismatch}: confidence and support statistics are
computed on the training distribution; the evaluation distribution may
differ, making confidence a poor predictor of action quality at evaluation
states.
\emph{Conflict-induced action suppression}: the confidence-ranked arbiter
always selects among the conflicting candidates; random selection may
occasionally select the action preferred by neither candidate's confidence
ordering but preferred by the environment, analogous to the beneficial
effect of exploration noise.
\emph{Action diversity}: Acrobot has three actions, and confidence-ranked
arbitration systematically favours one source agent's preferred action,
which may reduce diversity relative to random selection.
All three explanations are \emph{post hoc}; testing them requires
additional controlled experiments.

\paragraph{Steelman.}
The strongest objection to this work is that provenance-preserving
rule composition, in the form tested here, adds significant engineering
complexity for a modest and inconsistent performance benefit:
grammar fusion beats the best single agent by $42$--$57$ return units,
but loses substantially to the actor-logit ensemble on CartPole and
achieves the same as the ensemble on Acrobot, while requiring full
rule-bank induction, a shared symboliser, and an auditor.
A practitioner seeking audit capability without performance cost would
do better to log the neural ensemble's decisions with actor-index
attribution, which requires no symbolic machinery.
This objection has merit.
Our response is that the framework provides a qualitatively different
kind of auditability: it exposes \emph{which rules} were available,
\emph{what conflicts existed}, and \emph{which training records} back
each decision, in a form that an external auditor can verify without
access to the original model weights.
Whether that additional transparency justifies the engineering overhead
is a question for the application domain, not for this paper.

\paragraph{Implications.}
The results suggest two design recommendations for practitioners who
adopt rule-level composition.
First, the shared symboliser is a prerequisite for meaningful Jaccard
measurement; seed-specific symbolisers produce a Jaccard that partly
reflects quantiser-edge artefacts.
Second, the arbitration mechanism requires independent evaluation,
not just a pre-specified priority order; the Acrobot result shows that
a theoretically motivated ranking can be dominated by random selection
under distribution shift.

%% ════════════════════════════════════════════════════════════════════════════
\section{Limitations}
\label{sec:limitations}
%% ════════════════════════════════════════════════════════════════════════════

\paragraph{Experimental scale.}
The composition experiments use two deterministic benchmark environments,
300 training episodes per seed, and 100 evaluation episodes.
Performance on continuous-action tasks, stochastic environments, or
multi-step planning problems may differ substantially.

\paragraph{Arbitration mechanism.}
Confidence-ranked arbitration is not established as a generally good
conflict resolver; the Acrobot negative finding cannot be extrapolated
to other tasks without further experiments.
The mechanism underlying the random-rule advantage is untested.

\paragraph{Symbolic overlap as proxy.}
Jaccard under the shared symboliser is a syntactic statistic.
Experiment~B shows it overstates behavioral consensus in the tested
setting.
The degree of overstatement will depend on the task, the training-episode
budget, and the bin resolution.

\paragraph{Audit scope.}
The replay-and-provenance audit covers the logged execution on the tested
hardware and software configuration.
It does not cover the optimiser state, weight gradients, or random-number
sequences during training; training-process auditability is not established.
Cross-hardware or cross-software-version reproducibility is not claimed.

\paragraph{FQE quality.}
The fitted Q-evaluation baseline failed pre-specified quality checks
on both environments.
GPI-based composition may produce better results with a more powerful
critic architecture or more training data; the current negative result
applies to the specific FQE implementation and training budget used.

\paragraph{AIM scope.}
The causal-symbol result applies to a designed private-target coordination
game with a binary target signal.
Generalisation to downstream tasks, multi-step reasoning, or RL
control benchmarks requires separate experiments.

%% ════════════════════════════════════════════════════════════════════════════
\section{Conclusion}
\label{sec:conclusion}
%% ════════════════════════════════════════════════════════════════════════════

We presented a provenance-preserving rule-level merger for independently
trained reinforcement-learning agents and evaluated it empirically under
a shared frozen symboliser on two standard benchmarks with eight seeds each.

Three findings bear emphasis.
\emph{First}, exact environment replay and source-rule provenance
verification passed for all tested traces (${\approx}2.9$ million
environment transitions), establishing that the audit mechanism functions
as designed on the tested configuration.
\emph{Second}, the held-out behavioral agreement between the
highest-Jaccard seed pair in CartPole ($0.583$) was $51.4\%$, statistically
indistinguishable from the binary chance baseline, showing that symbolic
rule-set overlap is not a reliable proxy for neural-policy behavioral
consensus on fresh states.
\emph{Third}, confidence-ranked arbitration performed substantially worse
than random rule selection on Acrobot, indicating that arbitration design
is a primary performance driver that warrants independent empirical evaluation.

Grammar fusion achieved statistically significant return improvements
over the best single agent on both tasks, but the mechanism behind those
improvements differs between environments: in CartPole the rule core is
large and conflicts are rare; in Acrobot the rule core is near-empty and
the improvement comes despite (or because of) near-universal conflicts.

These results support bounded claims.
Provenance-linked symbolic rule composition makes selected trajectory
decisions auditable in the tested environments.
It does not establish general RL training transparency, does not
outperform neural ensembles, and does not provide a generally reliable
arbitration mechanism.
Future work should address arbitration mechanism design, return-weighted
rule induction, and extension to tasks where neural ensembles do not
saturate the return scale.

%% ════════════════════════════════════════════════════════════════════════════
%% Bibliography
%% ════════════════════════════════════════════════════════════════════════════
\bibliographystyle{plainnat}
\begin{thebibliography}{9}

\bibitem[Barreto et al.(2017)]{barreto2017successor}
Barreto, A., Dabney, W., Munos, R., Hunt, J.~J., Schaul, T.,
Silver, D., and van Hasselt, H. (2017).
Successor features for transfer in reinforcement learning.
In \emph{Advances in Neural Information Processing Systems (NeurIPS)}.

\bibitem[Barreto et al.(2017b)]{barreto2017gpi}
Barreto, A., Borsa, D., Quan, J., Schaul, T., and Silver, D. (2017).
Transfer in deep reinforcement learning using successor features and
generalised policy improvement.
In \emph{Proceedings of the International Conference on Machine Learning
(ICML)}, pages 501--510.

\bibitem[Lazaridou and Baroni(2020)]{lazaridou2018emergent}
Lazaridou, A. and Baroni, M. (2020).
Emergent multi-agent communication in the deep learning era.
\emph{arXiv:2006.02419}.

\bibitem[Lowe et al.(2019)]{lowe2019pitfalls}
Lowe, R., Foerster, J., Boureau, Y.-L., Pineau, J., and Dauphin, Y. (2019).
On the pitfalls of measuring emergent communication.
In \emph{Proceedings of AAMAS}.

\bibitem[Nagarajan et al.(2018)]{nagarajan2018deterministic}
Nagarajan, P., Warnell, G., and Stone, P. (2018).
Deterministic implementations for reproducibility in deep reinforcement
learning.
\emph{arXiv:1809.05676}.

\bibitem[Nevill-Manning and Witten(1997)]{nevill1997sequitur}
Nevill-Manning, C.~G. and Witten, I.~H. (1997).
Identifying hierarchical structure in sequences: A linear-time algorithm.
\emph{Journal of Artificial Intelligence Research}, 7:67--82.

\bibitem[Thakoor et al.(2022)]{thakoor2022geometric}
Thakoor, S., Rowland, M., Borsa, D., Dabney, W., Munos, R.,
and Barreto, A. (2022).
Generalised policy improvement with geometric policy composition.
In \emph{Proceedings of ICML}.

\bibitem[Verma et al.(2018)]{verma2018pirl}
Verma, A., Murali, V., Singh, R., Kohli, P., and Chaudhuri, S. (2018).
Programmatically interpretable reinforcement learning.
In \emph{Proceedings of ICML}, pages 5045--5054.

\end{thebibliography}

%% ════════════════════════════════════════════════════════════════════════════
\appendix
\section{Prediction Registry}
\label{app:predictions}

\begin{table}[h]
\centering
\caption{Pre-registered predictions and their status.}
\label{tab:predictions}
\begin{tabular}{llp{5cm}l}
\toprule
ID & Prediction & Evidence anchor & Status \\
\midrule
P1 & Grammar fusion $\ge$ best single on CartPole
   & Table~\ref{tab:composition}: $+42.3$ $[24.0, 61.9]$
   & \textbf{Supported} \\[4pt]
P2 & Grammar fusion $\ge$ best single on Acrobot
   & Table~\ref{tab:composition}: $+56.6$ $[39.3, 75.2]$
   & \textbf{Supported} \\[4pt]
P3 & Held-out behavioral agreement $> 85\%$
   & Experiment B: $51.4\%$ CI $[50.5\%, 52.3\%]$
   & \textbf{Contradicted} \\[4pt]
P4 & MDL ratio $< 1.0$ at all prefix lengths
   & Six prefixes: $0.513$--$0.575$
   & \textbf{Supported} \\
\bottomrule
\end{tabular}
\end{table}

\section{Per-Seed AIM Causal Intervention Results}
\label{app:aim}

\begin{table}[h]
\centering
\caption{Per-seed AIM matched-token causal effect
  $\Delta P(B \mid C)$ in the private-target coordination game.
  Active token position varied across seeds.}
\label{tab:aim}
\begin{tabular}{ccccc}
\toprule
Seed & Active position & $\Delta P(B \mid C)$ & Shuffled accuracy & MI (bits) \\
\midrule
0 & 1 & 0.9974 & 0.522 & 0.985 \\
1 & 1 & 0.9970 & 0.503 & 0.992 \\
2 & 0 & 0.9990 & 0.502 & 0.992 \\
3 & 1 & 0.9997 & 0.503 & 0.992 \\
4 & 0 & 0.9970 & 0.502 & 0.982 \\
5 & 1 & 0.9985 & 0.505 & 0.984 \\
\midrule
Mean & --- & 0.9981 & 0.506 & 0.988 \\
\bottomrule
\end{tabular}
\end{table}

\end{document}
